\label{chap:83-17}

\section{Archimedean and exceptional morphisms from a common enriched domain}
\label{p12-83-c17-s1:chap:doble}
\index{double projection}
\index{Monstrous Moonshine}

The double projection compares two evaluations of the common enriched domain
\(X_\infty^{\rm cal}\) already constructed in
\cref{p12-83-c17-s2:thm:lector-excepcional-global-rev6}. One produces
Archimedean values; the other, together with a calibration of its codomain, preserves
the combinatorial incidence that leads to codes, designs, and
automorphism groups. Both evaluations start from the unique initial state
\(X_\infty^{\rm cal}\) and preserve their common provenance in distinct
codomains.

\subsection{Domains of the \(2\) projections}

On \(X_\infty^{\rm cal}\), the cylinders contract to a point and the
boundary admits the incidence prolongation. In this residence, we fix
\(\mathbf{Arch}:=\mathbb R\). The Archimedean evaluation and the exceptional
evaluation are, respectively,
\[
\mathfrak R:X_\infty^{\rm cal}\longrightarrow\mathbf{Arch},
\qquad
\mathfrak E:X_\infty^{\rm cal}\times\mathscr C_{\rm exc}
\longrightarrow\mathbf{Exc}.
\]
Here \(\mathscr C_{\rm exc}\) contains the projective order, Paley gauge,
dodecad, incidence alignment, and lattice chamber
specified in each construction. The first map produces values; the
second produces codes, designs, lattices, and automorphisms. The preceding
global theorem is the source of compatibility and permits comparison of
invariants obtained through distinct operations without again intersecting
domains defined separately. The finite comparison in this work is
performed on the flag
\(B_K\subset H^-_\alpha\).

\begin{figure}[H]
\centering
\resizebox{\textwidth}{!}{\input{figures/fig_doble_proyeccion}}
\caption{Typed scheme of the two projections. The Archimedean evaluation
identifies states within the same fiber; the exceptional evaluation, relative
to its calibration, preserves incidence structure and automorphisms.}
\label{p12-83-c17-s1:fig:doble-proyeccion}
\end{figure}

\subsection{Topology of the space of histories}

The formulation by cylinders admits a standard realization as an
inverse limit. Let \(\mathsf S_n\) be the finite sets of states
admissible at depth \(n\), endowed with the discrete topology, and let
\[
 \tau_{n+1,n}:\mathsf S_{n+1}\longrightarrow\mathsf S_n
\]
the truncation maps. In the HMT realization these maps
are surjective: the prolongation law supplies a surviving continuation
of each prefix, and its projective compatibility is preserved in
the limit by
\cref{p12-83-c17-s2:thm:lector-excepcional-global-rev6}. This is the instance
of the system involved in the global double projection
\eqref{eq:frontal-doble-proyeccion}. We define
\[
 \mathcal X=\varprojlim_n\mathsf S_n.
\]
If \(x\ne y\), let \(m(x,y)\) be the first depth at which their components differ. The function
\[
 d(x,y)=3^{-m(x,y)},\qquad d(x,x)=0,
\]
is an ultrametric that induces the topology of the inverse limit.

\begin{theorem}[Space of compact histories]
\label{p12-83-c17-s1:thm:compacto-historias}
\label{thm:compacto-historias}
The space \(\mathcal X\) is compact and totally disconnected. When
every history admits bifurcations at arbitrarily large depths,
\(\mathcal X\) has no isolated points and is homeomorphic to Cantor
space.
\end{theorem}

\begin{proof}
The product \(\prod_n\mathsf S_n\) is compact. The equations
\(\tau_{n+1,n}(x_{n+1})=x_n\) define a closed subset, so that
\(\mathcal X\) is compact. The sets of histories with a fixed prefix
are simultaneously open and closed and form a basis; this proves total
disconnectedness. Under the hypothesis of arbitrarily deep branching,
no cylinder contains an isolated history. The characterization of
compact, metric, perfect, and totally disconnected spaces proves
the final assertion.
\end{proof}

The closed intervals \(I_n(s)\) associated by the Archimedean reader with
each \(s\in\mathsf S_n\) are compatible with truncation and satisfy
\[
 \operatorname{diam}I_n(s)\leq C\lambda^n,
 \qquad 0<\lambda<1.
\]
The intersection of the intervals corresponding to a history then contains a unique point.

\begin{theorem}[Archimedean Evaluation]
\label{p12-83-c17-s1:thm:evaluacion-holder}
\label{thm:evaluacion-holder}
The map
\[
 \operatorname{val}:\mathcal X\longrightarrow\mathbb R,
 \qquad
 \{\operatorname{val}(x)\}=\bigcap_n I_n(x_n),
\]
is continuous and Hölder with respect to \(d\). If
\(x\sim_{\mathrm{Arch}}y\) is equivalent to
\(\operatorname{val}(x)=\operatorname{val}(y)\), the induced application
is a homeomorphism
\[
 \mathcal X/{\sim_{\mathrm{Arch}}}
 \;\cong\;\operatorname{val}(\mathcal X).
\]
\end{theorem}

\begin{proof}
Two histories with a common prefix of length \(n\) have values in the
same interval of diameter no greater than \(C\lambda^n\). Since
\(d(x,y)=3^{-n}\), this estimate is a Hölder bound. The induced
map on the quotient is a continuous bijection from a compact space onto
a Hausdorff space and therefore a homeomorphism.
\end{proof}

The theorem formalizes the Archimedean projection: the value retains the evaluation
class and eliminates the distinction between histories situated in the same
fiber. Coherent coverage, compact exhaustion and the
exact reconstruction of the quotient are proved in
\cref{p12-83-c18-s1:thm:sobreyectividad-evaluacion-compacto,p12-83-c18-s1:prop:agotamiento-recta-real-u024,p12-83-c18-s2:hmtch:thm-cociente-real};
therefore, in the active HMT realization,
\(\operatorname{val}(\mathcal X)=\mathbb R\).

\subsection{Separation by the \(2\) projections}

\begin{theorem}[Joint separation]
\label{p12-83-c17-s1:thm:separacion-conjunta}
\label{thm:separacion-conjunta}
Let \(P_1:\mathcal X\to Y_1\) and \(P_2:\mathcal X\to Y_2\) be continuous maps
with Hausdorff codomains. If
\[
 x\ne y\quad\Longrightarrow\quad
 P_1(x)\ne P_1(y)\ \text{or}\ P_2(x)\ne P_2(y),
\]
the diagonal application
\[
 (P_1,P_2):\mathcal X\longrightarrow Y_1\times Y_2
\]
is a topological embedding.
\end{theorem}

\begin{proof}
The hypothesis expresses precisely the injectivity of the diagonal map. A
continuous injective map from a compact space to a Hausdorff space is a
homeomorphism onto its image.
\end{proof}

With a declared gauge \(c\in\mathscr C_{\rm exc}\) fixed, let
\[
 \mathfrak E_c:X_\infty^{\rm cal}\longrightarrow\mathbf{Exc},
 \qquad
 \mathfrak E_c(x):=\mathfrak E(x,c).
\]
Applied to HMT, the theorem takes
\[
 P_1=\mathfrak R,\qquad P_2=\mathfrak E_c
\]
on the common domain, or on the quotient by the declared coordinate
changes. Projective compatibility and the global double realization
are proved, respectively, in
\cref{p12-83-c17-s2:thm:lector-excepcional-global-rev6,p12-83-c17-s2:thm:doble-realizacion-global-rev6}.
Faithfulness is typed by the codomain: the exceptional reader is faithful on
the visible stream, while the enriched domain preserves the memory that
no isolated publication contracts. The passage to every depth is thus a
conclusion of the construction, not an external hypothesis.

At a finite level \(S\), this condition is decidable. Given
\(p:S\to A\), \(q:S\to B\), and an equivalence relation of coordinates,
the pairs \(x\not\sim y\) that simultaneously satisfy are enumerated
\[
 p(x)=p(y),\qquad q(x)=q(y).
\]
The absence of such pairs is equivalent to joint injectivity at that
level. The calculation must perform the two maps separately: when
\(q=f\circ p\), the second has the same fibers as the first and does not
recover information eliminated by \(p\).

At the level developed here, the equality
\[
 B_K=H_e\cap P_\pi
\]
constitutes the finite witness of compatibility between the two branches. Its
compatible prolongation to all depths is provided by
\cref{p12-83-c17-s2:thm:lector-excepcional-global-rev6,p12-83-c17-s2:thm:doble-realizacion-global-rev6};
the equality of supports is the finite shadow of that joint realization.

\subsection{Archimedean quotient by the evaluation fibers}

An HMT section contains a block, an orientation, a carry, an observable label,
and an extension coordinate. The
evaluation \(\val\) contracts its cylinders to a point. If two sections
\(x,y\) satisfy
\[
\val(x)=\val(y),
\]
the Archimedean evaluation identifies them although their transport is different. In
this precise sense,
\[
\boxed{\text{a real value is an equivalence class of the evaluation}.}
\]
On the canonical subdomain, sum, product and evaluation descend compatibly,
and the image inherits those operations. The existence of the limit, the
determination of its image, the characterization of the fibers and the coverage of
all of \(\RR\) are established by
\cref{p12-83-c18-s1:thm:sobreyectividad-evaluacion-compacto,p12-83-c18-s1:prop:agotamiento-recta-real-u024,p12-83-c18-s2:hmtch:thm-cociente-real}.
Finite windows exhibit the local structure of that global result; unitary
survival and compatibility are proved properties of the
HMT prolongation.

\subsection{Common invariant with the Witt system}

The vector \(K\) defines the support superior
\[
B_K=\{m:K_m\ge729\}=\{4,6,9,10\}.
\]
In the Witt atlas appear the supports
\[
P_\pi=\{2,4,5,6,7,8,9,10,11\},
\]
\[
H_e=\{1,3,4,6,9,10\},
\qquad
H_\varphi=\{1,2,3,4,7,9\}.
\]
The support of the negative carries of \(\alpha\) is
\[
H_\alpha^-=\{4,5,6,7,9,10\}.
\]
The support identities are verified.
\[
\boxed{
B_K=H_e\cap P_\pi=H_e\cap H_\alpha^-,}
\]
and
\[
\boxed{H_\alpha^-=B_K\sqcup\{5,7\}.}
\]
Witt polarity sends the four points of \(B_K\) to the pairs
\[
\{1,3\},\quad\{2,11\},\quad\{8,12\},\quad\{5,7\}.
\]
Three receive the positive label; the last completes the negative support.
These identities are exact in the declared gauge and constitute internal
checks of the same generated state: the vector, threshold, words and
Paley gauge are correlated coordinates of the construction, not
statistically independent inputs. The equality is obtained by
direct evaluation and does not constitute a probabilistic estimate.

\subsection{Residual design on a dodecada and elevation of hexads}
\label{p12-83-c17-s1:sec:w12-w24-elevacion}

The relation between the two Witt designs is obtained within the extended
binary Golay code and not through an identification of cardinalities.
Let
\[
\mathcal G_{24}\subset\mathbb F_2^{24}
\]
the extended binary Golay code, and denote by \(\mathcal O_{24}\) the family
of supports of its weight-eight words. It is a classical result that
\((\{1,\ldots,24\},\mathcal O_{24})\) is the Steiner system
\(W_{24}=S(5,8,24)\) \cite{MacWilliamsSloane1977,ConwaySloane1999}.
Fix a weight-twelve word and write \(D\) for its support. We define
\begin{equation}
 \mathcal H(D)
 :=\{O\cap D:O\in\mathcal O_{24},\ |O\cap D|=6\}.
 \label{p12-83-c17-s1:eq:hexadas-residuales-dodecada}
\end{equation}
The family \(\mathcal H(D)\) is intrinsic to the pair
\((\mathcal G_{24},D)\): it does not require choosing coordinates in \(D\).

\begin{theorem}[Residual design of a dodecada]
\label{p12-83-c17-s1:thm:dodecada-w12}
The incidence system
\[
 (D,\mathcal H(D))
\]
is a Steiner system \(S(5,6,12)\). In particular,
\(|\mathcal H(D)|=132\), and every subset of five points of \(D\) is
contained in a unique hexad of \(\mathcal H(D)\).
\end{theorem}

\begin{proof}
Let \(A\subset D\), with \(|A|=5\). Since the octads form
\(S(5,8,24)\), there is a unique octad \(O_A\) containing \(A\). If
\(j=|O_A\cap D|\), the binary sum of the characteristic words of
\(O_A\) and \(D\) belongs to \(\mathcal G_{24}\) and has weight
\[
 \operatorname{wt}(\mathbf1_{O_A}+\mathbf1_D)=8+12-2j=20-2j.
\]
The weights of the extended binary Golay code are
\(0,8,12,16,24\). Since \(j\leq8\), it follows that
\(j\in\{2,4,6\}\). The inclusion \(A\subset O_A\cap D\) gives \(j\geq5\),
hence \(j=6\). Therefore, \(O_A\cap D\in\mathcal H(D)\) is a hexad containing
\(A\).

If two members of \(\mathcal H(D)\) contained \(A\), their octads of
origin would contain the same set of five points; the Steiner property of
\(W_{24}\) would make them coincide. This proves uniqueness. The number
of blocks is obtained by counting incidences with five-point subsets:
\[
 |\mathcal H(D)|
 =\frac{\binom{12}{5}}{\binom65}=132.\qedhere
\]
\end{proof}

\begin{theorem}[Unique elevation of a hexad]
\label{p12-83-c17-s1:thm:elevacion-unica-hexada}
For each \(H\in\mathcal H(D)\), there exists a unique octad
\(O_H\in\mathcal O_{24}\) such that
\[
 O_H\cap D=H.
\]
Consequently, the application
\[
 \iota_D:\mathcal H(D)\longrightarrow\mathcal O_{24},
 \qquad H\longmapsto O_H,
\]
is well defined and is injective.
\end{theorem}

\begin{proof}
Existence forms part of the definition of \(\mathcal H(D)\). If two
octads \(O_1,O_2\) had intersection \(H\) with \(D\), both
would contain any five-point subset of \(H\). The uniqueness of
the octad incident with those five points implies \(O_1=O_2\).
Injectivity follows by intersecting again with \(D\).
\end{proof}

The lift has a local law that will be used to compare the five
regions of \(\pi\). In a system \(S(t,k,v)\), the number of blocks
containing a fixed subset of \(s\leq t\) points is
\[
 \lambda_s=\frac{\binom{v-s}{t-s}}{\binom{k-s}{t-s}}.
\]
Consequently,
\begin{equation}
 \lambda_4(W_{12})
 =\frac{\binom81}{\binom21}=4,
 \qquad
 \lambda_4(W_{24})
 =\frac{\binom{20}1}{\binom41}=5.
 \label{p12-83-c17-s1:eq:lambda4-w12-w24}
\end{equation}

\begin{corollary}[Decomposition \(4+1\) on a face]
\label{p12-83-c17-s1:cor:cuatro-mas-uno-w24}
Let \(B\subset D\), with \(|B|=4\). The five octads of \(W_{24}\) that
contain \(B\) decompose uniquely into
\[
 \{O_H:H\in\mathcal H(D),\ B\subset H\}
 \;\sqcup\;\{O_B^\perp\}.
\]
The first set contains four octads and satisfies
\(|O_H\cap D|=6\); the remaining octad satisfies
\[
 O_B^\perp\cap D=B.
\]
\end{corollary}

\begin{proof}
The first equality in \eqref{p12-83-c17-s1:eq:lambda4-w12-w24} provides four residual
hexads containing \(B\), and
\cref{p12-83-c17-s1:thm:elevacion-unica-hexada} provides four distinct octads. The
second equality in \eqref{p12-83-c17-s1:eq:lambda4-w12-w24} shows that there exists exactly
one fifth octad \(O_B^\perp\).

Since \(B\subset O_B^\perp\cap D\), the intersection spectrum used in
the proof of \cref{p12-83-c17-s1:thm:dodecada-w12} leaves only the possibilities four and
six. An intersection of size six would be a fifth residual hexad that
contains \(B\), contradicting \(\lambda_4(W_{12})=4\). Therefore, the
intersection has size four and coincides with \(B\).
\end{proof}

\subsubsection{Alineamiento with the \(12\) positions HMT}

The residual design is defined on the points of \(D\), whereas the
design obtained from the ternary code \(C_W\) is defined on the twelve
dodecaphase positions. To avoid conflating the two sets, we shall call a bijection
\emph{HMT binary alignment}
\begin{equation}
 \ell:\{1,\ldots,12\}_{\rm HMT}\xrightarrow{\sim}D
 \quad\text{such that}\quad
 \{\ell(H):H\in W_{12}^{\rm HMT}\}=\mathcal H(D).
 \label{p12-83-c17-s1:eq:alineamiento-hmt-w24}
\end{equation}
Here \(W_{12}^{\rm HMT}\) denotes the design of weight-six supports
obtained from the ternary code \(C_W\). Uniqueness of the Witt system
\(S(5,6,12)\) up to isomorphism ensures the existence of \(\ell\). Two
distinct alignments differ by automorphisms of the source and target
designs; hence \(\ell\) is coordinate data and not an additional
invariant.
The pair \((D,\ell)\) is the calibration datum required to transport
the lift \(\iota_D\) to the HMT\@ coordinate system. Once fixed, each
HMT hexad \(H\) unambiguously determines the octad
\[
 \widehat H:=\iota_D(\ell(H)).
\]
The dodecad and the alignment are data of the binary codomain; they are not identified
with a decimal output of the transport kernel.

For the distinguished face
\[
 B_K=\{4,6,9,10\},
\]
the four incident hexads are
\[
 B_K\sqcup\{1,3\},\quad
 B_K\sqcup\{2,11\},\quad
 B_K\sqcup\{5,7\},\quad
 B_K\sqcup\{8,12\}.
\]
The third agrees with
\(H_\alpha^-=B_K\sqcup\{5,7\}\). The alignment \(\ell\) lifts these
four hexads to the four residual octads over \(\ell(B_K)\);
\cref{p12-83-c17-s1:cor:cuatro-mas-uno-w24} also determines a unique transverse octad.

\subsubsection{Incidence correspondence between the microfiber of \texorpdfstring{\(\pi\)}{pi} and the residual octads: multiplicity \texorpdfstring{\(432+4\cdot144=1008\)}{432+4x144=1008}}

The exact microfiber of the word
\(w_\pi=010211\) contains five distinct decimal regions:
\[
\begin{array}{c|c|c}
\text{multiplicity}&U_6& e(U_6)\\
\hline
432&501|614|498|169|272|272&(5,5,5,5,5,5)\\
144&810|923|870|169|272|272&(3,3,9,5,5,5)\\
144&810|983|810|169|272|272&(3,9,3,5,5,5)\\
144&870|923|810|169|272|272&(9,3,3,5,5,5)\\
144&870|923|810|418|674|521&(9,3,3,7,1,7).
\end{array}
\]
Here \(e(U_6)\) is the multiplicative signature recovered from each
decimal block \(\overline{d_1d_2d_3}\) by means of
\(e=d_2-d_1\pmod {10}\). The first region is the only one of multiplicity
432 and constant signature. The remaining four have multiplicity 144; among
them, three share the return block \(169|272|272\), whereas the
fourth has return \(418|674|521\). The catalogue distinguishes the five
rows by a reproducible predicate: the three rows of type
\(\mathrm{ES}\) form the positive triple; among the two rows of type
\(\mathrm{NO}\), that of maximal multiplicity and constant signature is termed
transverse, and the remaining one is termed negative. Relative to this role
predicate, the partition is
\[
 \boxed{1_\perp+3_{++}+1_{--}.}
\]
The sign labels are a convention; what is intrinsic to the catalog is the
partition obtained from its columns, not a prior binary polarity.

The alineamiento regional complete the givesto \((D,\ell)\) with to bijection
that sends the transver region to \(O_{\ell(B_K)}^\perp\), the region
negativa to \(\widehat{H_\alpha^-}\) and the tris regionis positivas to the
elevacionis of
\[
 B_K\sqcup\{1,3\},\qquad
 B_K\sqcup\{2,11\},\qquad
 B_K\sqcup\{8,12\}.
\]
The ordered assignment between the three positive regions and those three hexads is a choice of coordinates. The two distinguished rows and the triple are reproducible through the preceding catalogue predicate; their designation as negative, transversal, and positive belongs to the alignment. Thus, the law \(4+1\) does not identify the five regions with five octads by mere cardinality: it provides the incidence space, whereas \((D,\ell)\) and the regional alignment specify the transport between the two objects.

\subsection{Lattice prolongation and dimensional screens}

The lift of hexads opens two branches that must not be confused. The first
conserves binary incidence and prolongs \(W_{12}\) toward \(W_{24}\). The
second conserves ternary transport:
\[
C_{12}^{(3)}
\longrightarrow N(A_2^{12})
\longrightarrow\Lambda_{24}
\longrightarrow V^\natural
\longrightarrow\mathbb M
\longrightarrow\text{Moonshine}.
\]
The primary proof of the gluing, the rootless \(3\) neighbor, and the corridor
to Moonshine resides in
\cref{chap:bandera-excepcional-acumulativa,thm:corredor-hmt-moonshine}.
There is no canonical arrow \(W_{24}\to N(A_2^{12})\), nor any direct
action of the Monster on TPK words: both branches start from the same
calibrated code and preserve different structures.

The reductions \(243/11\), \(12\to11\to10\), and \(26\to24\) are not
identities between numerals either. The \cref{chap:pantallas-dualidad} constructs their
spaces, quotients, and maps: the perfect ternary quotient, the screen
marked by \(K\), and the isotropic reduction of the Lorentzian lattice. This
section thereby retains a single function: to link the already proved finite
incidence with its two prolongations, without repeating the reference proofs.

\subsection{Residue–quotient involution and conservation of state information}

For \(n=9Q+r\), we define the lattice energy
\[
E_R(r,Q)=\frac{r^2}{R^2}+Q^2R^2.
\]
Then
\[
\boxed{E_R(r,Q)=E_{1/R}(Q,r).}
\]
The exchange \((r,Q,R)\leftrightarrow(Q,r,1/R)\) is an exact involution
of the lattice model: the residue and quotient exchange their functions.

\subsection{Finite compatibility and joint prolongation}

The finite witness \(B_K\), the carry pattern, and the Witt frame fix
local compatibility. Its prolongation is no longer left as an abstract
scheme: the \cref{p12-83-c17-s2:thm:lector-excepcional-global-rev6} constructs, block by
block, the exceptional reader with a unique transported initial gauge, and the
\cref{p12-83-c17-s2:thm:doble-realizacion-global-rev6} proves that this reading and the
Archimedean evaluation globalize over the same TPK section. Truncating the
history before reading is equivalent in both cases to truncating the output after
reading.

The Archimedean evaluation identifies histories by their limiting value. The
exceptional evaluation preserves the visible current and its Paley--Witt
incidence; it does not thereby recover the hidden memories that the visible emitter
does not publish. The double projection is therefore a joint construction with
two distinct codomains, not an identification between the continuum and the
exceptional branch.

\begin{remark}[Compatibility of the two projections]
The support \(B_K\), the carry pattern, and the Witt frame provide the
common invariants. The lattice chain prolongs this compatibility to a
rootless neighbor and, by classical classification, to the Leech lattice. The
binary branch toward \(W_{24}\) remains separate: its incidence morphism is
\(\iota_D\), defined after fixing the dodecad and alignment
\((D,\ell)\), and its codomain is neither a Niemeier lattice nor the
Leech lattice.
\end{remark}


\section{Exceptional projection of the enriched state: compatible incidence at all depths}
\label{p12-83-c17-s2:chap:lector-excepcional-global-rev6}
\index{exceptional reader}\index{double projection}
\index{Paley--Witt}\index{inverse system}

The finite witness of the double projection acquires global scope when the
exceptional reading is applied to every visible block and the gauge is
transported, rather than chosen anew, as the depth increases. This
section constructs that family. The domain remains the enriched TPK
history; the exceptional reading is faithful on its visible stream and does not
purport to replace the memory fields that this stream does not publish.

\subsection{The local operation on the 729 words}

In the Paley--Witt chart one fixes
\[
A_W=
\begin{pmatrix}
0&1&1&1&1&1\\
1&0&1&2&2&1\\
1&1&0&1&2&2\\
1&2&1&0&1&2\\
1&2&2&1&0&1\\
1&1&2&2&1&0
\end{pmatrix}
\in M_6(\mathbb F_3).
\]
Matrix multiplication gives
\[
A_W^{\mathsf T}=A_W,\qquad A_W^2=-I_6.
\]
For every conjugated operator \(A=fA_Wf^{-1}\), define
\[
\Gamma_A:\mathbb F_3^6\longrightarrow\mathbb F_3^{12},
\qquad
\Gamma_A(w)=(w,wA).
\]
Its image is a calibrated copy \(C_W(A)\) of the ternary Witt code.
Projection onto the first six coordinates is a left inverse:
\[
\operatorname{pr}_1\Gamma_A=\operatorname{id}_{\mathbb F_3^6}.
\]
Therefore, \(\Gamma_A\) is total and injective on the \(729\) words and
does not need to distinguish in advance \(\pi\), \(e\), \(\varphi\), or any
other output.

\subsection{Compatible extensions and uniqueness of the initial scale}

Let \(X_n\) be the set of admissible TPK prefixes of depth \(n\), with
truncations
\[
p_{n,m}:X_n\longrightarrow X_m,\qquad m\leq n.
\]
Fix an initial exceptional frame \(f_0\). Since the changes \(g_j\) form
part of the transport stored by the history, a calibrated prefix is the
pair \((x,f_0)\); we write
\[
 X_n^{\rm cal}:=X_n\times\{f_0\},\qquad
 p_{n,m}^{\rm cal}(x,f_0):=(p_{n,m}x,f_0),
\]
and
\[
 X_\infty^{\rm cal}
 :=\varprojlim_n(X_n^{\rm cal},p_{n,m}^{\rm cal}).
\]
In what follows we abbreviate \(p_{n,m}^{\rm cal}\) as \(p_{n,m}\).
The compatible visible emission is
\[
\varepsilon_n:X_n^{\rm cal}\longrightarrow(\mathbb F_3^6)^n,
\qquad
\varepsilon_n(x)=(w_0,\ldots,w_{n-1}),
\]
and satisfies
\[
\operatorname{pr}_{n,m}\varepsilon_n
=\varepsilon_m p_{n,m}.
\]
This application distinguishes its three types:
\[
|\mathcal U_6|=468,\qquad
|\operatorname{im}w_6|=243,\qquad
|\mathbb F_3^6|=729.
\]
The emission \(w_j\) is not identified with the complete state that conserves
leaf, quotient, carry, orientation, phase, and survival.

The exceptional gauge consists of a Witt chart, a dodecad, an
alignment, a lattice origin, and an oriented chamber. It is fixed only
once in the initial block. For
\(x\in X_n^{\rm cal}\), transport of the history supplies
\[
 g_j(x)\in\operatorname{GL}_6(\mathbb F_3),
 \qquad 0\le j<n-1,
\]
and \(g_j(x)\) depends only on the prefix
\(p_{n,j+1}(x)\). The frame is transported by
\[
 f_{j+1}(x)=g_j(x)f_j(x),\qquad
 A_j=f_jA_Wf_j^{-1}.
\]
Thus, the family \((A_j)\) is determined by the history and the initial gauge; it is not a sequence of independent choices.

Let \(\mathcal F_{\rm cal}\subseteq\operatorname{GL}_6(\mathbb F_3)\) be the
set of frames reachable from \(f_0\). The emission preserving the
transported gauge is
\[
 \widetilde\varepsilon_n:X_n^{\rm cal}
 \longrightarrow(\mathbb F_3^6\times\mathcal F_{\rm cal})^n,
 \qquad
 \widetilde\varepsilon_n(x)
 =\bigl((w_0,f_0),\ldots,(w_{n-1},f_{n-1})\bigr).
\]
Prefix dependence implies
\(\operatorname{pr}_{n,m}\widetilde\varepsilon_n
=\widetilde\varepsilon_m p_{n,m}\). Its projection onto the words is
\(\varepsilon_n\); forgetting the frames does not suffice to reconstruct the
exceptional reading.

\subsection{Naturality and compatibility at every depth}

The local operation is covariant under the row-vector convention. For
every change of coordinates \(f\in\operatorname{GL}_6(\mathbb F_3)\),
\[
A^f=fAf^{-1}
\quad\Longrightarrow\quad
\Gamma_{A^f}(wf^{-1})
=\Gamma_A(w)(f^{-1}\oplus f^{-1}).
\]
Change of frame simultaneously transports the word, the operator, and the
atlas. It is not required to remain within the stabilizer of a single
labeled copy.

With the set of reachable frames \(\mathcal F_{\rm cal}\) already defined,
we construct the fixed calibrated atlas
\[
 \mathscr C_W^{\rm cal}
 :=
 \left\{
 \bigl(f,\Gamma_{fA_Wf^{-1}}(w)\bigr):
 f\in\mathcal F_{\rm cal},\ w\in\mathbb F_3^6
 \right\}.
\]
The label \(f\) preserves the calibrated copy in which each word resides.
Define
\[
Q_n(x)=
\bigl(
\bigl(f_0,\Gamma_{A_0}(w_0)\bigr),\ldots,
\bigl(f_{n-1},\Gamma_{A_{n-1}}(w_{n-1})\bigr)
\bigr)\in\bigl(\mathscr C_W^{\rm cal}\bigr)^n.
\]
Let
\[
 r_{n,m}:
 \bigl(\mathscr C_W^{\rm cal}\bigr)^n
 \longrightarrow
 \bigl(\mathscr C_W^{\rm cal}\bigr)^m
\]
the truncation of the last \(n-m\) blocks.

\begin{theorem}[Projective compatibility of the exceptional incidence]
\label{p12-83-c17-s2:thm:lector-excepcional-global-rev6}
\label{thm:lector-excepcional-global-rev6}
For every \(m\leq n\),
\[
\boxed{r_{n,m}\circ Q_n=Q_m\circ p_{n,m}.}
\]
Consequently, there exists a unique application
\[
\boxed{
Q_\infty:X_\infty^{\rm cal}
\longrightarrow
\bigl(\mathscr C_W^{\rm cal}\bigr)^{\mathbb N}}
\]
whose finite projections are the \(Q_n\). The map is faithful on the
visible current.
\end{theorem}

\begin{proof}
Write
\[
\varepsilon_n(x)
=(w_0,\ldots,w_{m-1},w_m,\ldots,w_{n-1}).
\]
The compatibility of \(\varepsilon\) shows that the visible current of \(p_{n,m}(x)\) is \((w_0,\ldots,w_{m-1})\). For \(j<m\), the frame \(f_j\) depends only on the initial gauge and on changes preceding \(j\), all of which are contained in that prefix. The first \(m\) operators \(A_j\) are therefore the same before and after truncation. Applying \(\Gamma_{A_j}\) block by block yields the asserted identity. The universal property of the inverse limit produces \(Q_\infty\). Finally, \(\operatorname{pr}_1\Gamma_{A_j}=I\) recovers every \(w_j\), proving faithfulness on the visible current.
\end{proof}

The last assertion conserves its type. Two histories with the same current
\((w_j)_j\) and distinct hidden memory may have the same
exceptional image. The global reader preserves exactly what each visible block
publishes and the incidence added to that block.

The compatibility of enriched emissions induces a unique application.
\[
 \widetilde\varepsilon_\infty:X_\infty^{\rm cal}
 \longrightarrow
 (\mathbb F_3^6\times\mathcal F_{\rm cal})^{\mathbb N}.
\]
Let \(\operatorname{pr}_w\) the projection onto the corriente of palabras and
\[
 \widetilde\Gamma_\infty
 \bigl((w_j,f_j)_{j\ge0}\bigr)
 :=\bigl(
   (f_j,\Gamma_{f_jA_Wf_j^{-1}}(w_j))
 \bigr)_{j\ge0}.
\]
Then
\(Q_\infty=\widetilde\Gamma_\infty
\circ\widetilde\varepsilon_\infty\), with declared domains and codomains.

\subsection{Bifurcation of a surviving prefix in Archimedean value and exceptional incidence}

Fixed \(m\in\mathbb Z\), let us denote by
\[
 \operatorname{flat}:(\mathbb F_3^6)^{\mathbb N}
 \longrightarrow\mathbb F_3^{\mathbb N}
\]
the concatenation of the blocks in the declared coordinate order, and by
\[
 \operatorname{ev}_3((d_k)_{k\ge1})
 :=\sum_{k\ge1}\frac{d_k}{3^k}\in[0,1],
 \qquad
 \tau_m(t):=m+t,
\]
the ternary evaluation and translation by the integer part. The visible prefix
\((d_1,\ldots,d_{6n})\) determines the cylinder
\[
I_n(x)=
\left[
m+\sum_{k=1}^{6n}\frac{d_k}{3^k},
m+\sum_{k=1}^{6n}\frac{d_k}{3^k}+3^{-6n}
\right].
\]
A prolongation restricts the cylinder and
\(\operatorname{diam}I_n=3^{-6n}\to0\). Therefore, every compatible
infinite history determines a unique real point
\(\mathcal P_{\rm arq}(x)\).

\begin{theorem}[Double global realization]
\label{p12-83-c17-s2:thm:doble-realizacion-global-rev6}
\label{thm:doble-realizacion-global-rev6}
Every compatible section \(x=(x_n)_n\in X_\infty^{\rm cal}\) determines,
from the same prefixes, two compatible outputs:
\[
\mathcal P_{\rm arq}(x)\in\mathbb R,
\qquad
\mathcal P_{\rm exc}(x)=Q_\infty(x)
\in\bigl(\mathscr C_W^{\rm cal}\bigr)^{\mathbb N}.
\]
Truncating before reading is equivalent to truncating each reading after reading. The
first output identifies histories by their limiting value; the second preserves
each visible block and adjoins its Paley--Witt rotation.
\end{theorem}

\begin{proof}
The nesting of the cylinders and the convergence of their diameters prove the
first assertion. The
\cref{p12-83-c17-s2:thm:lector-excepcional-global-rev6} proves the second and the
compatibility with truncations. Both factor through the same emission:
\[
\mathcal P_{\rm arq}
=\tau_m\circ\operatorname{ev}_3\circ\operatorname{flat}
 \circ\operatorname{pr}_w
 \circ\widetilde\varepsilon_\infty,
\qquad
\mathcal P_{\rm exc}
=\widetilde\Gamma_\infty\circ\widetilde\varepsilon_\infty.
\]
\end{proof}

The double projection is thus constructed on admissible histories and a
transported initial gauge. The continuum is the Archimedean evaluation of
the section; the exceptional branch is the compatible realization of its
visible current. From \(C_W\), with the already declared alignments, the
branches toward \(W_{12}\), \(W_{24}\), and toward
\(N(A_2^{12})\), Leech, \(V^\natural\), and the Monster are opened. These branches do not turn
the digits into a set on which the Monster acts directly:
they transport the incidence that evaluation by value ceases to see.

\subsection{Postgenerative coupling with Euler's operator closure}
\label{p12-83-c17-s2:sec:acoplamiento-euler-doble-proyeccion-rev3}

Let \(x_\pi\in X_\infty^{\rm cal}\) be the surviving section of the
\(\pi\) channel. The double realization of
\cref{p12-83-c17-s2:thm:doble-realizacion-global-rev6} produces
\[
 \mathcal P_{\rm arq}(x_\pi)\in\mathbb R,
 \qquad
 \mathcal P_{\rm exc}(x_\pi)=Q_\infty(x_\pi).
\]
The first branch publishes the generated angle; the second preserves the visible
current and realizes the quadratic Paley--Witt relation in each block.

\begin{theorem}[Euler--postgenerative coupling double projection]
\label{p12-83-c17-s2:thm:acoplamiento-euler-doble-proyeccion-rev3}
Suppose that the coinductive construction of the channel of \(\pi\) determines
\[
 \mathcal P_{\rm arq}(x_\pi)=\pi_{\rm HMT},
\]
and that the exceptional branch uses the gauge transported with
\(A_W^2=-I_6\). Let \(J_{\mathrm e}\) be the real realization of the integral
source of \cref{thm:dos-realizaciones-fuente-cuadratica-rev3}. Then
\begin{equation}
 \boxed{
 \operatorname{Exp}\!\bigl(
   \mathcal P_{\rm arq}(x_\pi)J_{\mathrm e}
 \bigr)+I=0.}
 \label{p12-83-c17-s2:eq:acoplamiento-euler-doble-proyeccion-rev3}
\end{equation}
\end{theorem}

\begin{proof}
By the \cref{thm:dos-realizaciones-fuente-cuadratica-rev3}, \(A_W\) and
\(J_{\mathrm e}\) realizan the relation \(u^2+1=0\) in categorias distinct.
In the branch real,
\[
 \operatorname{Exp}(\theta J_{\mathrm e})
 =\cos\theta\,I+\sin\theta\,J_{\mathrm e}.
\]
Substitution of the previously generated output
\(\theta=\mathcal P_{\rm arq}(x_\pi)=\pi_{\rm HMT}\) produces \(-I\), and the
sum of \(I\) gives the stated identity.
\end{proof}

\paragraph{Causal order and absence of circularity.}
The demonstrated dependence is
\[
 \text{emitter}
 \longrightarrow x_\pi
 \longrightarrow\mathcal P_{\rm arq}(x_\pi)
 \longrightarrow\pi_{\rm HMT}
 \longrightarrow
 \operatorname{Exp}(\pi_{\rm HMT}J_{\mathrm e})+I.
\]
The exceptional branch fixes the quadratic type; it selects neither the
angle nor its digits. Euler's identity subsequently recognizes the
compatibility of the two realizations and does not enter as an input to the
generator of \(\pi_{\rm HMT}\).


\section{Natural globalization of the Archimedean and exceptional projections over a common inverse system}
\label{p12-83-c17-s3:chap:globalizacion-doble-proyeccion}

The double projection thesis acquires mathematical content when both
readings are defined on the same state system and commute with
depth restriction. Archimedean evaluation contracts a family of
cylinders to a value; exceptional evaluation preserves incidence
relations that the first may forget. These are not two juxtaposed
numerical expressions, but two natural transformations with a common
domain.

\subsection{Compatible systems at finite depth}

Let \((X_n,p_{n+1,n})\), \((A_n,a_{n+1,n})\) and
\((E_n,e_{n+1,n})\) be three inverse systems. The first contains TPK
histories; the second, Archimedean cylinders or approximants; the third, exceptional
incidence data. The active construction supplies the maps
\[
 R_n:X_n\longrightarrow A_n,
 \qquad
 Q_n:X_n\times C_n\longrightarrow E_n,
\]
where \(C_n\) collects the coordinate data of the exceptional reader. The
gauges form a compatible family and, for every \(n\), the diagrams
commute
\[
 a_{n+1,n}\circ R_{n+1}=R_n\circ p_{n+1,n},
\]
\[
 e_{n+1,n}\circ Q_{n+1}
 =Q_n\circ(p_{n+1,n}\times c_{n+1,n}).
\]

\begin{theorem}[Globalization of the two evaluations]
\label{p12-83-c17-s3:thm:globalizacion-evaluaciones}
By the preceding compatibilities, already constructed block by block, there exist
canonical maps
\[
 R_\infty:\varprojlim X_n\longrightarrow\varprojlim A_n,
 \qquad
 Q_\infty:\varprojlim X_n\times\varprojlim C_n
            \longrightarrow\varprojlim E_n.
\]
Its diagonal application
\[
 (R_\infty,Q_\infty):X_\infty\times C_\infty
 \longrightarrow A_\infty\times E_\infty
\]
is unique among the applications whose coordinate projections are
\(R_n\) and \(Q_n\).
\end{theorem}

\begin{proof}
Let \(x=(x_n)_n\in\varprojlim X_n\). The first commutativity equation
implies
\[
 a_{n+1,n}(R_{n+1}(x_{n+1}))=R_n(x_n),
\]
and therefore \((R_n(x_n))_n\) belongs to \(\varprojlim A_n\). We define
\(R_\infty(x)\) by means of that family. The same argument, applied to
\((x_n,c_n)_n\), defines \(Q_\infty\). Uniqueness follows because an element of
an inverse limit is determined by all of its coordinates.
\end{proof}

The theorem records the concrete proof carried out in HMT: the Archimedean
result and the exceptional code come from the same compatible
states and respect the truncation maps.

\subsection{Archimedean contraction of compatible cylinders and realization of the real line}

To each \(x_n\in X_n\) is associated a nonempty compact interval
\(I_n(x_n)\subset\mathbb R\). If
\[
 p_{n+1,n}(x_{n+1})=x_n
 \quad\Longrightarrow\quad
 I_{n+1}(x_{n+1})\subseteq I_n(x_n)
\]
and there exists \(\varepsilon_n\to0\) with
\(\operatorname{diam}I_n(x_n)\leq\varepsilon_n\), the intersection
\[
 \bigcap_{n\geq0}I_n(x_n)
\]
contains a unique real. Thus it is obtained.
\[
 \operatorname{ev}_{\rm arq}:X_\infty^{\rm arq}\longrightarrow\mathbb R.
\]
The phrase ``the cylinders tend to zero'' means here that their diameter
converges to zero, not that the evaluated number tends to zero.

The relation
\[
 x\sim_{\rm arq}y
 \quad\Longleftrightarrow\quad
 \operatorname{ev}_{\rm arq}(x)=\operatorname{ev}_{\rm arq}(y)
\]
describes exactly the information removed by the reader. The real value
is an evaluation class; the history additionally preserves orientation,
carry, sheet, and boundary signature.

\subsection{Exceptional projection by incidence and boundary conservation}

The second evaluation does not take values in \(\mathbb R\). Its codomain is a
category of incidence configurations. In the current finite realization,
the coordinate datum contains:
\[
 \begin{aligned}
 C_{\rm exc}=(\,&\text{projective order},\text{Paley gauge},
 \text{dodecad},\\
 &\text{alignment},\text{lattice origin},\text{chamber}).
 \end{aligned}
\]
After fixing that data, the chain contains two descendants that must remain separated:
\[
 C_W\longrightarrow W_{12}\longrightarrow W_{24},
\]
\[
 C_W\longrightarrow N(A_2^{12})
 \longrightarrow\Lambda_{24}
 \longrightarrow V^\natural
 \longrightarrow\mathbb M.
\]
The first transports designs; the second leads, through later
classical results, to the Leech lattice, the vertex operator algebra, and
the Monster group. The central holonomy and the
Schläfli--\(E_6\) corridor form a parallel boundary extraction; they are not obtained
by replacing one of the two preceding descents.

\subsection{Finite Compatibility of the Common Domain}

The dodecaphonic seal
\[
 K=(234,543,140,729,659,824,621,58,914,794,146,601)
\]
determines the support
\[
 B_K=\{m:K_m\geq729\}=\{4,6,9,10\}.
\]
In the published Witt chart,
\[
 P_\pi=\{2,4,5,6,7,8,9,10,11\},
 \qquad
 H_e=\{1,3,4,6,9,10\},
\]
and, by independent evaluation of the two supports,
\[
 \boxed{B_K=H_e\cap P_\pi.}
\]
Furthermore, for the support of negative carry
\[
 H_\alpha^-=\{4,5,6,7,9,10\}
\]
we have
\[
 B_K=H_e\cap H_\alpha^-,
 \qquad H_\alpha^-=B_K\sqcup\{5,7\}.
\]
These identities constitute a witness of finite compatibility: the same
subset is recovered from the arithmetic seal and from Witt
incidence. This is not an equality of cardinalities, since it identifies specific
positions.

\subsection{Typed separation of the Archimedean and exceptional publications of the common state}

\begin{theorem}[Joint Fidelity Criterion]
Let \(X\) be compact and \(Y,Z\) be Hausdorff spaces. If two
continuous maps \(R:X\to Y\) and \(Q:X\to Z\) satisfy
\[
 x\neq y\Longrightarrow R(x)\neq R(y)\ \text{or}\ Q(x)\neq Q(y),
\]
then \((R,Q):X\to Y\times Z\) is a topological embedding.
\end{theorem}

\begin{proof}
The condition is the injectivity of the diagonal application. Every continuous
injective application from a compact space into a Hausdorff space is a
homeomorphism onto its image.
\end{proof}

At each finite depth, the condition is decided by enumerating the
pairs \(x\neq y\) that remain simultaneously in a fiber of \(R_n\) and
of \(Q_n\). The compatibility of those checks when passing from \(n+1\) to \(n\)
is precisely the projective identity proved in
\cref{p12-83-c17-s2:thm:lector-excepcional-global-rev6}, and
\cref{p12-83-c17-s2:thm:doble-realizacion-global-rev6} carries out its
globalization. The witness \(B_K\) exhibits the common finite invariant; the
obstruction of the minimal reader in
\cref{p12-83-c17-s4:thm:obstruccion-lector-excepcional-minimo} delimits the
visible quotient, while the enriched reader preserves the orientation,
sheet, signature, return and boundary data required by the double
publication.

\subsection{Precise formulation of the thesis on the real line}

The enriched object associated to the Archimedean reading is the arrow
\[
 \mathbb R_{\mathfrak H}:=
 \bigl(X_\infty^{\rm arq}
 \xrightarrow{\operatorname{ev}_{\rm arq}}
 \operatorname{Im}(\operatorname{ev}_{\rm arq})\subseteq\mathbb R\bigr).
\]
The domain preserves genealogies; the image preserves values. Coverage of
the cylinders at every scale and compatible prolongability are proved in
the construction of the Archimedean reader
\eqref{eq:continuo-lector-arquimediano}: for every compact set
\(K\subset\mathbb R\),
\(\operatorname{ev}_{\mathbb R}(\Omega_{\infty,K})=K\), and the directed union
of those intervals covers \(\mathbb R\). Therefore,
\[
 \operatorname{Im}(\operatorname{ev}_{\rm arq})=\mathbb R,
 \qquad
 X_\infty^{\rm arq}/\!\sim_{\mathbb R}\ \cong\mathbb R.
\]
The real field does not act as the starting point of the construction, but as
the Archimedean quotient of a space of sections with memory. The descent of
sum and product is verified operation by operation and does not condition this
already proved surjectivity.

The double reading is then expressed by
\[
 X_\infty^{\rm com}\times C_\infty
 \xrightarrow{\ (\operatorname{ev}_{\rm arq},Q_\infty)\ }
 \mathbb R\times E_\infty.
\]
This formula is the mathematical translation of the conceptual assertion: one
and the same history may contract to a value and, in the transverse reading
of its boundary, preserve the incidence leading to the exceptional
symmetries.


% Fuente normativa activa de la obstrucción de factorización.
\section{Dependence of the reader of incidence with respect to the enriched state}
\label{p12-83-c17-s4:sec:obstruccion-lector-excepcional-minimo}

Finite compatibility of the two evaluations does not imply that every
incidence reader refines the partition produced by the visible reader. The
regional catalog makes it possible to decide this question exactly for the
minimal Paley--Witt realization. Let \(\mathcal U_{468}\) be the set of
the 468 published regional states, and let
\[
 \Psi:\mathcal U_{468}\longrightarrow\mathcal W_{243}\subset\mathbb F_3^6,
 \qquad
 \Psi(U_1,\ldots,U_6)=(-U_1,\ldots,-U_6)\pmod 3.
\]
For \(w\in\mathbb F_3^6\), the map
\[
 c(w)=(w,wA_W)\in C_W\subset\mathbb F_3^{12}
\]
is the ternary encoder fixed in the Paley--Witt section. We define
the minimal incidence reader by
\[
 Q_{\min}=\operatorname{Supp}\circ c\circ\Psi:
 \mathcal U_{468}\longrightarrow\mathcal P([12]).
\]

For a map \(f:X\to Y\), denote by
\[
 \operatorname{Eq}(f)=\{(x,x')\in X^2:f(x)=f(x')\}
\]
its fiberwise equivalence relation.

\begin{theorem}[Obstruction of the minimal exceptional reader]
\label{p12-83-c17-s4:thm:obstruccion-lector-excepcional-minimo}
In the complete catalogue of 468 states one verifies
\[
 \operatorname{Eq}(\Psi,Q_{\min})=\operatorname{Eq}(\Psi).
\]
In the usual abbreviated notation for fiber partitions,
\[
 \ker(\Psi,Q_{\min})=\ker\Psi.
\]
Consequently, \((\Psi,Q_{\min})\) is not injective and the minimal
incidence reader does not distinguish any pair of states that has already been
identified by \(\Psi\).
\end{theorem}

\begin{proof}
The factorization \(Q_{\min}=S\circ\Psi\), with
\(S=\operatorname{Supp}\circ c\), implies
\[
 \Psi(u)=\Psi(v)\Longrightarrow Q_{\min}(u)=Q_{\min}(v).
\]
The converse implication between the equivalence relations is immediate,
because \(\Psi\) is the first component of the diagonal map. Equality of the
relations therefore follows without statistical assumptions. Exhaustive
enumeration further confirms that both partitions have the same
histogram:
\[
 132\,[1]+66\,[2]+18\,[3]+3\,[4]+6\,[5]+18\,[6],
\]
where \(a\,[m]\) means that there exist \(a\) fibres of cardinality \(m\).
\end{proof}

The microfiber of \(w_\pi=010211\) provides the distinguished witness. Its
five regions are
\[
\begin{split}
 &(501,614,498,169,272,272),\quad
 (810,923,870,169,272,272),\\
 &(810,983,810,169,272,272),\quad
 (870,923,810,169,272,272),\\
 &(870,923,810,418,674,521).
\end{split}
\]
The five have the same visible image and the same exceptional support.
\[
 P_\pi=\{2,4,5,6,7,8,9,10,11\}.
\]
At the level of the 243 words there appear 213 distinct supports: 185 have
one preimage, 27 have two, and one has four. Taking support therefore introduces
a second quotient in addition to the one imposed by \(\Psi\).

The theorem delimits the minimal reader, not the enriched exceptional
evaluation. A reader capable of separating the regions of the same word
must act before the visible quotient and use at least one of the data
that it eliminates: orientation, sheet, signature, return, or boundary state. The
dodecaphase intertwiner developed in the following part belongs to that
enriched class and does not contradict the preceding obstruction.

Complete enumeration of the fibers produces the preceding histogram and
proves that taking support introduces a second quotient.


\section{Archimedean and incidental publication}

For each level, let \(\mathcal E_n\) be an incidence object and let
\begin{equation}
 e_n:S_n\longrightarrow\mathcal E_n,
 \qquad
 \sigma_{n,m}:\mathcal E_n\longrightarrow\mathcal E_m
 \label{p12-83-c17-s5:eq:datos-lector-incidencial-u024}
\end{equation}
applications that satisfy
\begin{equation}
 \sigma_{n,m}\circ e_n=e_m\circ r_{n,m}.
 \label{p12-83-c17-s5:eq:naturalidad-lector-incidencial-u024}
\end{equation}
The universal property constructs
\begin{equation}
 \operatorname{ev}_{\rm inc}:\Omega_\infty
 \longrightarrow\mathcal E_\infty,
 \qquad
 \mathcal E_\infty:=\varprojlim_n\mathcal E_n.
 \label{p12-83-c17-s5:eq:evaluacion-incidencial-limite-u024}
\end{equation}
The double publication is
\begin{equation}
 \mathcal P_\infty
 :=(\operatorname{ev}_{\mathbb R},\operatorname{ev}_{\rm inc}):
 \Omega_\infty\longrightarrow\mathbb R\times\mathcal E_\infty.
 \label{p12-83-c17-s5:eq:doble-publicacion-continuo-u024}
\end{equation}

\begin{figure}[H]
\centering
\begin{tikzcd}[column sep=10em,row sep=5em,
  cells={nodes={font=\large}}]
&\Omega_\infty
 \arrow[dl,"\operatorname{ev}_{\mathbb R}"']
 \arrow[dr,"\operatorname{ev}_{\rm inc}"]&\\
\mathbb R
 \arrow[dr,"\operatorname{coord}"']&&
\mathcal E_\infty
 \arrow[dl,"\operatorname{pub}_{\rm inc}"]\\
&\mathcal O&
\end{tikzcd}
\caption{Two publications of the same history. Commutativity in
\(\mathcal O\) does not identify the codomains \(\mathbb R\) and
\(\mathcal E_\infty\).}
\label{p12-83-c17-s5:fig:doble-publicacion-continuo-u024}
\end{figure}

The naturality proved in
\cref{p12-83-c17-s2:thm:lector-excepcional-global-rev6,p12-83-c17-s2:thm:doble-realizacion-global-rev6}
is expressed, after applying the publication readers, by
\begin{equation}
 \operatorname{pub}_{\rm inc}\circ\operatorname{ev}_{\rm inc}
 =\operatorname{coord}\circ\operatorname{ev}_{\mathbb R},
 \label{p12-83-c17-s5:eq:compatibilidad-doble-publicacion-u024}
\end{equation}
and produces a typed comparison of the two publications. The identity
preserves the common origin without constructing one branch from the other: both
start from \(\Omega_\infty\).


\section{Gauge of the exceptional publication}

The second projection is not a unary reader devoid of incidence data. Let
\[
 \mathscr C_{\rm exc}
 =\mathscr C_{\rm ord}\times\mathscr C_{\rm Paley}
  \times\mathscr C_{\rm dod}\times\mathscr C_{\rm inc}
  \times\mathscr C_{\rm ret}
\]
the space of exceptional gauges, whose coordinates respectively fix
the projective order, the Paley gauge, the dodecad, the incidence
alignment, and the lattice chamber. One fixes
\[
 \chi_{\rm exc}
 =(\chi_{\rm ord},\chi_{\rm Paley},\chi_{\rm dod},
   \chi_{\rm inc},\chi_{\rm ret})\in\mathscr C_{\rm exc}.
\]
The typed reader has domain
\[
 \mathfrak E:
 X_\infty^{\rm cal}\times\mathscr C_{\rm exc}
 \longrightarrow\mathbf{Exc},
 \qquad
 \mathfrak E_{\chi_{\rm exc}}(u):=\mathfrak E(u,\chi_{\rm exc}).
\]
Henceforth, every abbreviation \(\mathfrak E(u)\) means
\(\mathfrak E_{\chi_{\rm exc}}(u)\) for this explicitly fixed gauge.
The dependence on \(\chi_{\rm exc}\) does not affect the existence of the
continuum or its Archimedean publication: it types the exceptional conservation
of incidence once the common state has been constructed.

\begin{falsifier}[Loss of the exceptional gauge]
The comparison becomes incomplete if the projective order, the dodecad,
the incidence alignment, or the lattice chamber is changed without transporting the remaining
gauge coordinates, or if \(\mathfrak E\) is used as a unary reader before
fixing \(\chi_{\rm exc}\). Such a failure invalidates that exceptional comparison;
it does not modify the Archimedean quotient or regenerate the continuum.
\end{falsifier}


The second reading does not arise from the real number. Let
\[
 \mathfrak E:
 \mathscr O_{\mathbb R,\leftrightarrow}^{\rm HMT}
 \longrightarrow\mathbf{Exc}
\]
the continuous incidence reader, with Hausdorff codomain, that prolongs
the Paley--Witt, Golay,
Leech, \(V^\natural\), Monster, and Moonshine chain. Both maps begin from
the same history:
\begin{equation}
 \mathfrak D(u)=(\mathfrak R(u),\mathfrak E(u)).
 \label{p12-83-c17-s7:hmtch:eq-doble-proyeccion-two-way}
\end{equation}
The first contracts the compatible fibres to their Archimedean mean; the
second preserves incidence, orientation, and boundary memory.

Fixing the origin of the dodecaphonic calendar, let it be
\[
 R_+:=\{k\in\mathbb Z:1+(k\bmod 9)\in\{1,4,7\}\}
\]
the bi-infinite set of complete returns \((\tau=+1)\). It is syndetic:
its gaps are uniformly bounded. For
\(u=(m,(x_k)_{k\in\mathbb Z})\), define
\[
 \operatorname{Ret}_+(u)
 :=\bigl(m,(x_k)_{k\in R_+}\bigr)
 \in\mathbb Z\times\prod_{k\in R_+}X,
\]
but each \(x_k\) here is the whole state
\eqref{p12-83-c19-s1:hmtch:eq-celula-trit-completa} with leaf, carry, route, and incidence,
not the isolated sign.

\begin{proposition}[Lossless Sampling of the Return Genealogical Record]
\label{p12-83-c17-s7:hmtch:prop-retorno-reconstruye}
The map \(\operatorname{Ret}_+\), with the subspace topology in the
product, is a homeomorphism onto its image.
Consequently, there is a unique continuous map
\(\mathfrak E_+\) on that image such that
\begin{equation}
 \mathfrak E=\mathfrak E_+\circ\operatorname{Ret}_+.
 \label{p12-83-c17-s7:hmtch:eq-factorizacion-excepcional-retorno}
\end{equation}
\end{proposition}

\begin{proof}
Between two consecutive returns there is a uniformly bounded number of steps
of the calendar. The complete state at a return contains the transition and
the genealogical register that determine the intermediate steps; in the other direction, the
inverse arrow of the groupoid completion is used. The entire
orbit is thereby reconstructed. The map is continuous and, when restricted to a
compact chart, a continuous injection into a Hausdorff space is a
homeomorphism onto its image. The integer coordinate is preserved
explicitly. The factorization is defined by transporting \(\mathfrak E\) through
the inverse and is continuous chart by chart.
\end{proof}

For this reason, \(+1\) participates in the exceptional corridor as a complete memory
of return: sampling the integer states does not lose the exceptional output.
This proposition proves recovery and factorization, not a generation of
the corridor from the isolated scalar \(+1\). The Archimedean output and the
exceptional output are two coordinated resultants of the same antecedent.
